\paragraph{C.3ah (Refuted, with explicit counterexamples): the regular ambient range $N=r+2$.}
For every integer $r\ge1$, there is a regular excellent strictly
henselian local complex ring of dimension $N=r+2$ and a principal
divisor with regular puncture such that
\[
 \mathrm{CH}^r(X)=0,
 \qquad
 \mathrm{CH}^r(Y)_{\mathrm{alg}}\otimes\mathbb Q\ne0.
\]
In particular ordinary $J^r$ restriction is not surjective,
even with rational coefficients. For $r=2$ this gives the
dimension-four case not covered by C.3ae or C.3ag.

\begin{proof}
Put $d=r+2$ and specify
\[
 \begin{aligned}
 R&=\mathbb C[x_0,\ldots,x_{r+1}]_{(x_0,\ldots,x_{r+1})}^{\mathrm{sh}},
 & f&=x_0^d+\cdots+x_{r+1}^d,\\
 X&=\operatorname{Spec}R\setminus\{\mathfrak m_R\},
 &Y&=X\cap V(f).
 \end{aligned}
\]
Strict henselization is taken at the evident complex geometric
point. The hypersurface cone has dimension $r+1\ge2$.
Its gradient is $(d x_i^{d-1})_i$, so its only singular point
is the vertex. It is integral and normal: its projective
base is smooth and connected, and the hypersurface is
Cohen--Macaulay and regular in codimension one. The quotient
$R/(f)$ is the strict henselization of its vertex local ring;
this is the usual compatibility of strict henselization with
quotients, with the same geometric residue point. Thus $R$
and $Y$ satisfy all the required local hypotheses.

We use the smooth projective base
\[
 T=V(x_0^d+\cdots+x_{r+1}^d)\subset\mathbb P^{r+1},
 \qquad H=c_1(\mathcal O_T(1)),
 \qquad D=T\cap V(x_{r+1}).
\]
Let $k_0=\overline{\mathbb Q}\subset\mathbb C$, choose
$u_1,\ldots,u_r$ algebraically independent over $k_0$, and choose
$w,\eta$ with
\[
 w^d=-1-\sum_{i=1}^r u_i^d,
 \qquad \eta^d=-1.
\]
We take the specified points and degree-zero cycle
\[
 P=(u_1:\cdots:u_r:w:1),\qquad
 a=(1:\eta:0:\cdots:0)\in D(k_0),\qquad
 \alpha=[P]-[a].
\]
The point $P$ is a geometric generic point of the
$k_0$-variety $T$. The class $\alpha$ is ordinarily
algebraically trivial: when $r=1$ the curve $T$ itself
provides the difference family; when $r\ge2$, a smooth
connected complete-intersection curve through $P,a$
provides it. Such a curve is obtained by Bertini with the
two prescribed smooth points, in a sufficiently ample
linear system. The diagonal on that curve followed by its
closed immersion in $T$ gives an actual fixed-product
cycle family.

We recall why this particular class survives the hyperplane
quotient, rather than merely appealing to nonvanishing of
some zero-cycle. Suppose a nonzero multiple of $\alpha$
belonged to $H\mathrm{CH}^{r-1}(T)_{\mathbb Q}$. Every
hyperplane product has a representative supported on $D$.
After clearing denominators, the finitely many cycles and
rational functions expressing the supposed relation can be
spread, specialized to $\overline{k_0(T)}$, and descended
to a finite extension of $k_0(T)$. Norm and closure then
give
\[
       b\Delta_T=\Theta_D+\Theta_Z,\qquad b\ne0,
\]
where $\Theta_D$ is supported on $T\times D$ and
$\Theta_Z$ on $Z\times T$ for a proper divisor $Z\subset T$.
These finite-data and localization operations are those
proved in C.3f and C.3s. Adjunction gives $K_T=\mathcal O_T$,
so $H^{r,0}(T)\ne0$. Both correspondences on the right
annihilate this space: restriction to a resolution of $Z$
kills an $r$-form, and pushforward from $D$ cannot produce
Hodge bidegree $(r,0)$. The diagonal acts as multiplication
by $b$, a contradiction. Consequently
\[
 \alpha\ne0\quad\hbox{in}\quad
 \frac{\mathrm{CH}^r(T)_{\mathbb Q}}
      {H\mathrm{CH}^{r-1}(T)_{\mathbb Q}}.
\]

The Fermat hypersurface is projectively normal with section
ring generated in degree one. The strict-local cone injection
C.3e therefore sends $\alpha$ to a non-torsion class
$\gamma\in\mathrm{CH}^r(Y)$. Its construction is the flat
pullback from $T$ to the punctured line bundle, followed by
the strict-local base change. Pulling back the fixed curve
family just described proves
$\gamma\in\mathrm{CH}^r(Y)_{\mathrm{alg}}$ integrally.
This is exactly the case $m=0$ of C.3s.

Finally the local map $\mathbb C\to R$ is geometrically
regular, Noetherian, and has henselian target. We apply
\citet[Corollary~6.4(i), author p.~27]{rootSkalit2017}
with zero-dimensional base. It gives
$\mathrm{CH}^r(\operatorname{Spec}R)=0$ for every $r>0$,
integrally. Chow localization surjects this group onto
$\mathrm{CH}^r(X)$, so the source is zero. The specified
non-torsion class $\gamma$ proves the failure of ordinary
restriction surjectivity.
\end{proof}

Combining $r=2$ here with C.3r at $N=5$, C.3ag at
$N=6$, and C.3ae at every $N\ge7$ yields the exact
universal dimension statement over $\mathbb C$:
\[
 \text{ordinary }J^2(X)\longrightarrow J^2(Y)
 \text{ is universally surjective}
 \quad\Longleftrightarrow\quad N\le3.
\]
For $N\le3$ the target has no codimension-two points and
its Chow group is zero. For every $N\ge4$ the cited
examples satisfy the regular-puncture hypotheses and give
failure after tensoring with $\mathbb Q$. The polynomial
and dimension checks for the present construction are
retained in \texttt{code/check\_regular\_ambient\_calabi\_yau.py}.
Its checks support the displayed equations; the generic-point
and Hodge-form proof above establishes non-torsion in all
dimensions.
